\scanpage{89}
\mspage{41}

17. Récapitulation sur les $\mathcal V_i$, $\mathcal V_{i_0\ldots i_n}$, $\mathcal V_J$, $\mathcal V_J^L$, $\mathcal V_{J,L}$, et introduction des $\mathcal V_{J_0,J_1\ldots J_n}$.

Soit $I$ un ens. ordonné. Une partie \inserted{$J$} de $I$ est dite \uline{fermée} si
\[j\in J,\quad i\leq j\quad\Longrightarrow\quad i\in J.\]
Les parties fermées de $I$ forment une \deleted{\unsure{ens.}} partie de $\mathcal P(I)$ stable par intersections quelconques et par réunions \deleted{finies}, quelconques, ce sont donc \deleted{les} \inserted{les} fermés d'une top. de $I$, appelée sa \uline{top.} \inserted{inférieure} \deleted{\uline{canonique}. inf. Soit $L=I\setminus J$, \unsure{on a}}. On a donc pour toute partie $A$ de $I$
\[
(82)\qquad\bar A=\{i\in I\mid\exists j\in A\text{ tel que }i\leq j\}
\]
\margin{(adh. de $A$ pour la top. \unsure{can.} de $I$)}

\unsure{Soit} $\unsure{A|J}\subset I$, $L=I\setminus\unsure{\bar A|\bar J}$, pour que $J$ soit fermé, \inserted{i.e. $L$ soit \uline{ouverte},} il f. et s. que $L$ soit ``\uline{cofermé}'' i.e. fermé pour la top. associée à l'ordre opposé à \deleted{\unsure{I}} celui de $I$, i.e.
\[l\in L,\quad\deleted{\unsure{l\in I}}\ i\geq l\quad\Longrightarrow\quad i\in L.\]
L'adhérence d'une partie $B$ de $I$ pour la top. supérieure est notée $\widetilde B$ :
\[
(83)\qquad\widetilde B=\{i\in I\mid\exists j\in B\text{ tel que }i\geq j\}
\]
\deleted{\unsure{Donc}} Ainsi $\widetilde B$ est aussi le plus petit ouvert (pour la top. inf.) contenant $B$.

\inserted{\uline{Pour}} Une partie $J$ de $I$, il revient au même qu'elle soit \unsure{loc\textsuperscript{t}|loc!} fermée (i.e. intersection d'un fermé et d'un ouvert), ou qu'on ait $J=\bar J\cap\widetilde J$, ou encore
\margin{88}
